\documentclass{article}
\usepackage[utf8]{inputenc}
\usepackage{amsmath,amssymb}
\usepackage[most]{tcolorbox}
\usepackage{geometry}
\geometry{margin=2.5cm}

\newcommand{\step}[1]{\par\noindent\hangindent=3.5em\hangafter=1%
  \makebox[3.5em][l]{\textbf{#1.}}}
\newcommand{\steptag}[1]{\hfill{\small\textsf{[#1]}}}

\newtcolorbox{proofbox}[1]{
  colback=gray!5, colframe=gray!60,
  fonttitle=\bfseries, title={#1},
  breakable, enhanced,
  left=4pt, right=4pt, top=4pt, bottom=4pt,
  before skip=10pt, after skip=10pt
}

\begin{document}

\begin{proofbox}{Single-compression transfer: there are universal C\_co < i...}
\step{1} Single-compression transfer: there are universal C\_co < infinity and e\_co > 0 such that restricting an extended alpha-inclusion to one ideal and following it by one compatible amplified compression produces an extended C\_co*(alpha+epsilon)-inclusion whenever alpha+epsilon <= e\_co. \steptag{validated} \\[6pt]
\step{1.1} Precise compatible setup and universal ledger: let A be an extended epsilon-C*-algebra, B an exact unital C*-algebra, J a unital direct-summand ideal with unit q, and v:B->A an extended alpha-inclusion. The compatible single compression is defined by P=v(q) and T=Co\_P composed with v restricted to J. For e=alpha+epsilon below one universal threshold, P is a nonvanishing alpha-projection, S\_P is an extended C\_ca*e-C*-algebra, and universal constants D,K<infinity used below may be chosen independently of B,J,A,n. \steptag{validated} \\[6pt]
\step{1.1.1} Projection induced by the ideal unit. Since q=q\textasciicircum{}dagger=q\textasciicircum{}2 and ||q||=1 in the exact C*-algebra B, def-extended-delta-inclusion gives P\textasciicircum{}dagger=v(q\textasciicircum{}dagger)=P, ||P\textasciicircum{}2-P||=||v(q)v(q)-v(q\textasciicircum{}2)||<=alpha, and 1-alpha<=||P||<=1+alpha. Thus P is an alpha-projection by def-delta-projection and is nonvanishing for small e because its norm lies in the second alternative with abs(||P||-1)<=alpha<=e. \steptag{validated} \\[4pt]
\step{1.1.2} Target algebra and constant ledger. Let C\_ca,e\_ca be the universal witnesses in lem-compcb-corner-algebra. Applied with delta=alpha, that external makes S\_P, with compressed product, inherited involution and u\_P=Co\_P(P), an extended C\_ca*e-C*-algebra whenever e<=e\_ca. By the operator comparison asserted in def-compressed-corner there are universal k\_cmp<infinity and e\_cmp0>0 such that at every amplified level ||Co\_\{P\_n\}(Z)-(P\_n Z)P\_n||<=k\_cmp*e*||Z||. Let D=2*k\_cmp+5 and K\_N=D+1. Together with the universal witnesses C\_r,e\_r from lem-compcb-rectangular-product and the thresholds from lem-compcb-amplified-compression and lem-compcb-amplified-compression-identities, all constants and their finite minimum are universal. \steptag{validated} \\[4pt]
\step{1.2} At every n>=1, T\_n=id\_\{M\_n\} tensor T equals Co\_\{P\_n\} composed with v\_n on M\_n tensor J, has range in M\_n tensor S\_P, is linear, preserves dagger exactly, and sends the source unit I\_n tensor q exactly to the amplified compressed unit Co\_\{P\_n\}(P\_n). \steptag{validated} \\[6pt]
\step{1.2.1} Amplified formula and range. By lem-compcb-amplified-compression, id\_\{M\_n\} tensor Co\_P=Co\_\{P\_n\} and M\_n tensor S\_P=S\_\{P\_n\}, where P\_n=I\_n tensor P. Therefore T\_n(x)=Co\_\{P\_n\}(v\_n(x)) for x in M\_n tensor J and lies in Image(Co\_\{P\_n\})=S\_\{P\_n\}=M\_n tensor S\_P. Both v\_n and Co\_\{P\_n\} are linear by def-extended-delta-inclusion and def-compressed-corner, so T\_n is linear. \steptag{validated} \\[4pt]
\step{1.2.2} Exact dagger and unit clauses. The star clause of the extended alpha-inclusion gives v\_n(x\textasciicircum{}dagger)=v\_n(x)\textasciicircum{}dagger. Lem-compcb-amplified-compression-identities with the square pair P\_n,P\_n gives Co\_\{P\_n\}(Z\textasciicircum{}dagger)=Co\_\{P\_n\}(Z)\textasciicircum{}dagger, hence T\_n(x\textasciicircum{}dagger)=T\_n(x)\textasciicircum{}dagger. The unit of M\_n tensor J is q\_n=I\_n tensor q and v\_n(q\_n)=I\_n tensor v(q)=P\_n, so T\_n(q\_n)=Co\_\{P\_n\}(P\_n)=I\_n tensor Co\_P(P)=I\_n tensor u\_P, where the last equality again uses lem-compcb-amplified-compression. Thus the unit defect is zero. \steptag{validated} \\[4pt]
\step{1.2.3} Context-dependent derivation. Import validated node 1.1: J has unit q, P=v(q), T=Co\_P composed with v restricted to J, and P is an alpha-projection, with e=alpha+epsilon below the universal thresholds for the amplified-compression results. For n>=1 put v\_n=id\_\{M\_n\} tensor v, P\_n=I\_n tensor P, and q\_n=I\_n tensor q. Tensor functoriality for linear maps gives T\_n=id\_\{M\_n\} tensor (Co\_P composed with v|\_J)=(id\_\{M\_n\} tensor Co\_P) composed with v\_n. Lem-compcb-amplified-compression changes this exactly to Co\_\{P\_n\} composed with v\_n and identifies Image(Co\_\{P\_n\})=S\_\{P\_n\}=M\_n tensor S\_P, proving the formula and range; linearity follows because v\_n and Co\_\{P\_n\} are linear. For x in M\_n tensor J, the exact star clause for the extended alpha-inclusion and lem-compcb-amplified-compression-identities for the square pair (P\_n,P\_n) give T\_n(x\textasciicircum{}dagger)=Co\_\{P\_n\}(v\_n(x)\textasciicircum{}dagger)=Co\_\{P\_n\}(v\_n(x))\textasciicircum{}dagger=T\_n(x)\textasciicircum{}dagger. Finally v\_n(q\_n)=I\_n tensor v(q)=P\_n, so T\_n(q\_n)=Co\_\{P\_n\}(P\_n)=I\_n tensor Co\_P(P), the last equality being lem-compcb-amplified-compression. Thus every clause of node 1.2 follows in the setup supplied by node 1.1. \steptag{validated} \\[4pt]
\step{1.3} Uniform compression-closeness estimate: there are universal D<infinity and e\_d>0 such that, whenever e<=e\_d, every n>=1 and x in M\_n tensor J satisfy ||T\_n(x)-v\_n(x)||<=D*e*||x||. \steptag{validated} \\[6pt]
\step{1.3.1} Ideal-support estimates before compression. Put q\_n=I\_n tensor q, P\_n=v\_n(q\_n), and Z=v\_n(x). Since q\_n x=x q\_n=x in M\_n tensor J, the multiplicative clause of the extended alpha-inclusion gives ||P\_n Z-Z||<=alpha||x|| and ||ZP\_n-Z||<=alpha||x||. Its norm bounds also give ||P\_n||<=1+alpha and ||Z||<=(1+alpha)||x||. \steptag{validated} \\[4pt]
\step{1.3.2} Compression comparison and assembly. Def-compressed-corner, applied in the epsilon-C*-algebra M\_n tensor A to the alpha-projection P\_n, gives ||Co\_\{P\_n\}(Z)-(P\_n Z)P\_n||<=k\_cmp*e||Z|| with the universal k\_cmp of node 1.1.2. If e<=1/2, the epsilon-C*-product bound and node 1.3.1 yield ||(P\_n Z)P\_n-Z||<=||(P\_n Z-Z)P\_n||+||ZP\_n-Z||<=[(1+epsilon)(1+alpha)+1]*alpha||x||<=5e||x||. Also k\_cmp*e||Z||<=2*k\_cmp*e||x||. Thus ||T\_n(x)-v\_n(x)||<=D*e||x|| for D=2*k\_cmp+5, uniformly in n. \steptag{validated} \\[4pt]
\step{1.3.3} The validated external lem-compcb-amplified-almost-containment is consistent with and strengthens the same compatibility principle when a restricted range is already presented as an inner compressed corner: under its stated projection-containment hypotheses it directly supplies a universal amplified compression-closeness bound. The present canonical ideal compression needs no extra range hypothesis because nodes 1.3.1-1.3.2 derive that bound directly from q\_n x=x q\_n=x. \steptag{validated} \\[4pt]
\step{1.4} The compression-closeness estimate and the two-sided bounds for the extended alpha-inclusion imply (1-K\_N*e)||x||<=||T\_n(x)||<=(1+K\_N*e)||x|| at every n, for the universal K\_N=D+1. \steptag{validated} \\[6pt]
\step{1.4.1} For any n and x, node 1.3 gives ||T\_n(x)-v\_n(x)||<=D*e||x||, while def-extended-delta-inclusion gives (1-alpha)||x||<=||v\_n(x)||<=(1+alpha)||x||. The triangle and reverse-triangle inequalities therefore give ||T\_n(x)||<=[1+alpha+D*e]||x||<=[1+(D+1)e]||x|| and ||T\_n(x)||>=[1-alpha-D*e]||x||>=[1-(D+1)e]||x||. This is uniform in n and proves the claim with K\_N=D+1. \steptag{validated} \\[4pt]
\step{1.5} With the compressed product on S\_P, there is a universal K\_M<infinity such that every n and x,y in M\_n tensor J satisfy ||T\_n(xy)-T\_n(x) dot T\_n(y)||<=K\_M*e*||x||||y|| whenever e is below the common threshold. \steptag{validated} \\[6pt]
\step{1.5.1} Four-term multiplicativity telescope. For fixed n and x,y in M\_n tensor J, insert v\_n(xy), v\_n(x)v\_n(y), and T\_n(x)T\_n(y) to obtain ||T\_n(xy)-T\_n(x) dot T\_n(y)|| <= ||T\_n(xy)-v\_n(xy)|| + ||v\_n(xy)-v\_n(x)v\_n(y)|| + ||v\_n(x)v\_n(y)-T\_n(x)T\_n(y)|| + ||T\_n(x)T\_n(y)-T\_n(x) dot T\_n(y)||. This uses only the ambient and compressed products already defined. \steptag{validated} \\[4pt]
\step{1.5.2} Uniform bounds for the four terms. The first is at most D*e||x||||y|| by node 1.3 and exact source submultiplicativity; the second is at most alpha||x||||y|| by the extended inclusion. For the third, insert T\_n(x)v\_n(y), use the epsilon-C*-product bound, node 1.3, ||v\_n(z)||<=(1+alpha)||z||, and node 1.4 to get at most 2*D*(K\_N+3)*e||x||||y|| when e<=1. Since T\_n(x),T\_n(y) lie in the amplified square corner, lem-compcb-rectangular-product gives for the fourth at most C\_r*e||T\_n(x)||||T\_n(y)||<=C\_r*(1+K\_N)\textasciicircum{}2*e||x||||y||. Therefore K\_M:=D+1+2*D*(K\_N+3)+C\_r*(1+K\_N)\textasciicircum{}2 is finite, universal, and bounds the sum. \steptag{validated} \\[6pt]
\step{1.5.2.1} Bounds for the first two telescope terms, conditional on the registered compression estimate. By node 1.3 applied to xy, ||T\_n(xy)-v\_n(xy)||<=D*e||xy||<=D*e||x||||y||, since M\_n tensor J is an exact C*-algebra and hence its product is submultiplicative. By the multiplicative clause of the extended alpha-inclusion, ||v\_n(xy)-v\_n(x)v\_n(y)||<=alpha||x||||y||<=e||x||||y||. \steptag{validated} \\[6pt]
\step{1.5.2.1.1} Source-algebra typing and norm control from the validated setup. Node 1.1 says that B is an exact unital C*-algebra and J is a unital direct-summand ideal of B. Hence J, with the inherited multiplication, involution, and norm, is a C*-subalgebra: for x,y in M\_n tensor J, the matrix product xy again lies in M\_n tensor J. The standard matrix amplification M\_n tensor J is therefore an exact C*-algebra, so its C*-norm is submultiplicative and ||xy||<=||x||||y||. Node 1.1 also supplies e=alpha+epsilon with alpha,epsilon>=0 and v:B->A an extended alpha-inclusion. \steptag{validated} \\[4pt]
\step{1.5.2.1.2} Apply validated node 1.3 to the now correctly typed element xy in M\_n tensor J: ||T\_n(xy)-v\_n(xy)||<=D*e||xy||<=D*e||x||||y||. Since v is an extended alpha-inclusion, def-extended-delta-inclusion says v\_n is an alpha-homomorphism at every n; its multiplicative clause gives ||v\_n(xy)-v\_n(x)v\_n(y)||<=alpha||x||||y||<=e||x||||y|| because alpha<=alpha+epsilon=e. These are precisely the first two telescope bounds. \steptag{validated} \\[4pt]
\step{1.5.2.2} Bound for the third telescope term, conditional on the registered compression and norm estimates. Insert T\_n(x)v\_n(y). The epsilon-C*-product bound, node 1.3, the inclusion norm bound, and node 1.4 give ||v\_n(x)v\_n(y)-T\_n(x)T\_n(y)|| <= (1+epsilon)[||v\_n(x)-T\_n(x)|| ||v\_n(y)||+||T\_n(x)|| ||v\_n(y)-T\_n(y)||] <= D*e*(1+epsilon)*[(1+alpha)+(1+K\_N*e)]||x||||y||. Since 0<=alpha,epsilon<=e<=1, this is at most 2*D*(K\_N+3)*e||x||||y||. \steptag{validated} \\[4pt]
\step{1.5.2.3} Product naturality and the fourth-term bound. Node 1.2.1 gives T\_n(x),T\_n(y) in S\_\{P\_n\}=M\_n tensor S\_P. If A=[a\_ij] and B=[b\_ij] lie there, the level-n product obtained by amplifying the compressed product on S\_P has entries sum\_k Co\_P(a\_ik b\_kj), which by linearity equal Co\_P(sum\_k a\_ik b\_kj). Thus A dot\_n B=(id\_\{M\_n\} tensor Co\_P)(AB)=Co\_\{P\_n\}(AB), using lem-compcb-amplified-compression. It is therefore exactly the compatible square-corner product estimated by lem-compcb-rectangular-product, not a different product. Consequently ||T\_n(x)T\_n(y)-T\_n(x) dot\_n T\_n(y)||<=C\_r*e||T\_n(x)|| ||T\_n(y)||<=C\_r*(1+K\_N)\textasciicircum{}2*e||x||||y||, where node 1.4 and e<=1 give the last inequality. \steptag{validated} \\[4pt]
\step{1.5.2.4} Summing the four bounds from nodes 1.5.2.1--1.5.2.3 in the telescope of validated node 1.5.1 gives [D+1+2*D*(K\_N+3)+C\_r*(1+K\_N)\textasciicircum{}2]*e||x||||y||. Hence K\_M:=D+1+2*D*(K\_N+3)+C\_r*(1+K\_N)\textasciicircum{}2 is finite and universal. The explicit dependency and requires-validated edges ensure this conclusion cannot be accepted before nodes 1.3 and 1.4; the range fact is taken from already validated node 1.2.1, and product naturality was proved in node 1.5.2.3, so no dependency on the circular sibling 1.5.3 is needed. \steptag{validated} \\[4pt]
\step{1.5.3} The rectangular-product application is legitimate at every n: for each z in M\_n tensor J, T\_n(z)=Co\_\{P\_n\}(v\_n(z)) and idempotence of Co\_\{P\_n\} puts T\_n(z) in S\_\{P\_n\}=M\_n tensor S\_P. Hence T\_n(x),T\_n(y) form the compatible amplified square-corner pair (P\_n,P\_n,P\_n); product naturality identifies the amplified S\_P product with Co\_\{P\_n\}(T\_n(x)T\_n(y)), so lem-compcb-rectangular-product applies with a threshold independent of n. \steptag{validated} \\[6pt]
\step{1.5.3.1} Product naturality under amplification: for every n and A=(a\_ij),B=(b\_ij) in M\_n tensor S\_P, if dot\_amp denotes the multiplication on M\_n tensor S\_P obtained by amplifying the compressed product a dot\_P b=Co\_P(ab), then (A dot\_amp B)\_ij=sum\_k Co\_P(a\_ik b\_kj)=Co\_P(sum\_k a\_ik b\_kj). Therefore A dot\_amp B=(id\_\{M\_n\} tensor Co\_P)(AB)=Co\_\{P\_n\}(AB), where the final equality is lem-compcb-amplified-compression. By def-compressed-corner, Co\_\{P\_n\}(AB) is exactly the rectangular compressed product for the compatible triple (P\_n,P\_n,P\_n), not merely an element of the same underlying space. \steptag{validated} \\[4pt]
\step{1.5.3.2} Local range proof, independent of nodes 1.2 and 1.5.2. Fix n>=1 and z in M\_n tensor J. By validated node 1.1, T=Co\_P composed with v|\_J, P=v(q), and P is an alpha-projection. Hence T\_n(z)=(id\_\{M\_n\} tensor Co\_P)(v\_n(z))=Co\_\{P\_n\}(v\_n(z)), where P\_n=I\_n tensor P, by lem-compcb-amplified-compression. By lem-compcb-amplified-compression-identities, Co\_\{P\_n\} is idempotent, so Co\_\{P\_n\}(T\_n(z))=Co\_\{P\_n\}\textasciicircum{}2(v\_n(z))=T\_n(z). Thus T\_n(z) belongs to Img(Co\_\{P\_n\})=S\_\{P\_n\} by def-compressed-corner; and lem-compcb-amplified-compression gives S\_\{P\_n\}=M\_n tensor S\_P. Applying this to z=x and z=y proves that both actual factors T\_n(x),T\_n(y) lie in M\_n tensor S\_P. All invoked thresholds depend only on the universal externals and not on n. \steptag{validated} \\[4pt]
\step{1.5.3.3} Now apply validated node 1.5.3.1 to A=T\_n(x), B=T\_n(y), whose hypotheses are supplied by the preceding local range proof: the amplified compressed-corner product equals Co\_\{P\_n\}(T\_n(x)T\_n(y)), exactly the rectangular compressed product for the compatible triple (P\_n,P\_n,P\_n). Therefore lem-compcb-rectangular-product applies and yields ||T\_n(x) dot\_amp T\_n(y)-T\_n(x)T\_n(y)|| <= C\_co*e*||T\_n(x)||*||T\_n(y)||, with a universal threshold independent of n. This establishes the legitimacy assertion without using pending node 1.2 or pending node 1.5.2. \steptag{validated} \\[4pt]
\step{1.6} Taking C\_co=max\{C\_ca,K\_N,K\_M\} and e\_co to be the positive minimum of all thresholds used above proves that T is an extended C\_co*(alpha+epsilon)-inclusion into the compressed corner; the constants are universal and no amplification level or ideal dimension enters. \steptag{validated} \\[6pt]
\step{1.6.1} Let e\_co be the minimum of 1/2, e\_ca, e\_cmp0, e\_r, e\_d and the positive thresholds in lem-compcb-amplified-compression and lem-compcb-amplified-compression-identities, and put C\_co=max\{1,C\_ca,K\_N,K\_M\}. Node 1.1 gives the target compressed corner as an extended C\_ca*e-C*-algebra. Node 1.2 proves at every n that T\_n is linear, has the correct target, preserves dagger, and has zero unit defect. The validation-gated direct estimate in node 1.6.1.2 gives multiplicative defect at most K\_M*e, without using pending node 1.5, and node 1.4 gives the two-sided (1 plus-or-minus K\_N*e) norm bounds. Consequently the guarded assembly in node 1.6.1.4 applies def-extended-delta-inclusion to show that every amplification is a C\_co*e-inclusion. Since e=alpha+epsilon and all constants and thresholds are universal, T is the required extended C\_co*(alpha+epsilon)-inclusion. \steptag{validated} \\[6pt]
\step{1.6.1.1} Dependency-explicit final assembly. After nodes 1.1, 1.2, 1.4, and 1.5 have been validated, fix n>=1. Node 1.1 supplies the target extended C\_ca*e-C*-algebra and universal constants/thresholds; node 1.2 supplies linearity, target membership, exact dagger preservation, and zero unit defect for T\_n; node 1.5 supplies multiplicative defect at most K\_M*e*||x||||y||; and node 1.4 supplies (1-K\_N*e)||x||<=||T\_n(x)||<=(1+K\_N*e)||x||. For C\_co=max\{1,C\_ca,K\_N,K\_M\}, these are respectively bounded by the C\_co*e target-algebra parameter, C\_co*e-homomorphism defects, and C\_co*e two-sided norm distortion. Hence T\_n is a C\_co*e-inclusion by def-extended-delta-inclusion. Since n was arbitrary, T is an extended C\_co*e-inclusion. With e=alpha+epsilon and e\_co the stated positive finite minimum of the universal thresholds, this is exactly the claimed extended C\_co*(alpha+epsilon)-inclusion, with constants independent of n and the ideal dimension. \steptag{validated} \\[4pt]
\step{1.6.1.2} Direct multiplicativity estimate without using pending node 1.5. Fix n>=1 and x,y in M\_n tensor J. The validated telescope 1.5.1 bounds the defect by four terms. For the first, validated 1.1 identifies J as an ideal in the exact C*-algebra B, hence M\_n tensor J is an exact C*-algebra and ||xy||<=||x||||y||; validated 1.3 applied to xy gives at most D*e||x||||y||. For the second, def-extended-delta-inclusion applied to the extended alpha-inclusion v gives at most alpha||x||||y||<=e||x||||y||. Validated 1.5.2.2 bounds the third term by 2*D*(K\_N+3)*e||x||||y||, and validated 1.5.2.3 bounds the fourth by C\_r*(1+K\_N)\textasciicircum{}2*e||x||||y||. Therefore, with K\_M:=D+1+2*D*(K\_N+3)+C\_r*(1+K\_N)\textasciicircum{}2, one has ||T\_n(xy)-T\_n(x) dot T\_n(y)||<=K\_M*e||x||||y||. All constants and thresholds are universal by 1.1, and the estimate is uniform in n. \steptag{validated} \\[4pt]
\step{1.6.1.3} Dependency-correct final assembly, independent of pending node 1.5. Once direct multiplicativity node 1.6.1.2 and validated nodes 1.1, 1.2, and 1.4 are validated premises, fix n>=1. Node 1.1 supplies the target extended C\_ca*e-C*-algebra and the universal threshold ledger; node 1.2 supplies linearity, correct target, exact dagger preservation, and zero unit defect for T\_n; node 1.6.1.2 supplies multiplicative defect at most K\_M*e||x||||y||; and node 1.4 supplies (1-K\_N*e)||x||<=||T\_n(x)||<=(1+K\_N*e)||x||. With C\_co=max\{1,C\_ca,K\_N,K\_M\}, every required defect and target parameter is at most C\_co*e, so T\_n is a C\_co*e-inclusion by def-extended-delta-inclusion. Since n is arbitrary, T is an extended C\_co*e-inclusion. Substituting e=alpha+epsilon and taking e\_co to be the positive minimum of the universal thresholds gives the root contract, uniformly in n and the ideal dimension. This derivation has no dependency on node 1.5. \steptag{archived} \\[4pt]
\step{1.6.1.4} Guarded final assembly replacing the unsound dependency on pending node 1.5. Require validation of nodes 1.1, 1.2, 1.4, and 1.6.1.2. Then for arbitrary n, node 1.1 gives the extended C\_ca*e-C*-algebra target and universal thresholds, node 1.2 gives the linear, target, dagger, and zero-unit-defect clauses, node 1.6.1.2 gives multiplicative defect <=K\_M*e||x||||y||, and node 1.4 gives the two-sided (1 plus-or-minus K\_N*e) norm bounds. For C\_co=max\{1,C\_ca,K\_N,K\_M\}, def-extended-delta-inclusion makes T\_n a C\_co*e-inclusion. Universality and arbitrariness of n make T an extended C\_co*e-inclusion; e=alpha+epsilon and the universal threshold minimum yield the contract. No premise from pending node 1.5 is used. \steptag{validated} \\[4pt]
\end{proofbox}

\end{document}
